\documentclass[11pt,a4paper]{article}

% ── Packages (only those available on this installation) ──────────────────────
\usepackage[margin=2.5cm]{geometry}
\usepackage{amsmath,amssymb}
\usepackage{graphicx}
\usepackage{booktabs}
\usepackage{longtable}
\usepackage{array}
\usepackage{xcolor}
\usepackage{hyperref}
\usepackage{caption}
\usepackage{subcaption}
\usepackage{enumitem}
\usepackage{microtype}
\usepackage{listings}
\usepackage{float}
\usepackage{placeins}
\IfFileExists{todonotes.sty}{\usepackage{todonotes}}{\newcommand{\todo}[2][]{}}

% ── Figure path ───────────────────────────────────────────────────────────────
\graphicspath{{figures/}}

% ── Hyperref setup ────────────────────────────────────────────────────────────
\hypersetup{
  colorlinks = true,
  linkcolor  = blue!70!black,
  citecolor  = green!50!black,
  urlcolor   = blue!70!black,
}

% ── Custom commands ───────────────────────────────────────────────────────────
\newcommand{\numu}{$\nu_\mu$}
\newcommand{\numubar}{$\bar{\nu}_\mu$}
\newcommand{\ccpip}{$\mathrm{CC}\pi^{\pm}$}
\newcommand{\pmu}{$p_\mu$}
\newcommand{\ppi}{$p_\pi$}
\newcommand{\costhmu}{$\cos\theta_\mu$}
\newcommand{\costhpi}{$\cos\theta_\pi$}
\newcommand{\thmupi}{$\theta_{\mu\pi}$}
\newcommand{\uboone}{MicroBooNE}
\newcommand{\GeVc}{\,\text{GeV}\!/c}
\newcommand{\pot}{\,\text{POT}}
\newcommand{\cm}{\,\text{cm}}

% inline math shorthand so they work inside math mode too
\newcommand{\pmuM}{p_\mu}
\newcommand{\ppiM}{p_\pi}
\newcommand{\cthmuM}{\cos\theta_\mu}
\newcommand{\cthpiM}{\cos\theta_\pi}
\newcommand{\thmupM}{\theta_{\mu\pi}}

% ── Listings (C++ snippets) ───────────────────────────────────────────────────
\lstset{
  language=C++,
  basicstyle=\ttfamily\small,
  keywordstyle=\bfseries\color{blue!70!black},
  commentstyle=\itshape\color{gray},
  backgroundcolor=\color{gray!10},
  frame=single,
  framerule=0.4pt,
  numbers=none,
  breaklines=true,
}

% ── Table column helpers ──────────────────────────────────────────────────────
\newcolumntype{C}[1]{>{\centering\arraybackslash}p{#1}}
\newcolumntype{R}[1]{>{\raggedleft\arraybackslash}p{#1}}
\newcolumntype{L}[1]{>{\raggedright\arraybackslash}p{#1}}

% Let TeX stretch spacing on hard-to-break lines (long \texttt paths/code) to keep
% text out of the margin.
\emergencystretch=3em

% ─────────────────────────────────────────────────────────────────────────────


% ---- cross-document references to the technical supplement ----
\usepackage{xr}
\externaldocument[S-]{technical_supplement}
\newcommand{\suppref}[2]{#1~\ref{S-#2} of the technical supplement}
\externaldocument[CR-]{cr_data_note}
\externaldocument[N-]{analysis_note}
\newcommand{\noteref}[2]{#1~\ref{N-#2} of the inclusive analysis note}
\newcommand{\crref}[2]{#1~\ref{CR-#2} of the control-sample note}
\usepackage{tikz}
\usetikzlibrary{arrows.meta,positioning}
% Every result figure shows fake data. The stamp is drawn over the graphic by LaTeX so it cannot be
% dropped by regenerating the figure, and so black points are never mistaken for beam-on data.
\newcommand{\blindfig}[2][\textwidth]{%
  \begin{tikzpicture}
    \node[inner sep=0pt] (img) {\includegraphics[width=#1]{#2}};
    \node[anchor=north east, fill=red!12, draw=red!60!black, rounded corners=1pt, inner sep=2pt,
          font=\scriptsize\bfseries\sffamily, text=red!60!black] at ([xshift=-2pt,yshift=-2pt]img.north east)
          {BLINDED: CV PSEUDO-DATA};
  \end{tikzpicture}}


\begin{document}
\begin{titlepage}
  \centering
  \vspace*{2cm}
  {\LARGE\bfseries Measurement of Proton-Tagged Charged-Current Single-Charged-Pion Production on
    Argon in the NuMI Beam with the \uboone{} Detector\par}
  \vspace{0.8cm}
  {\large Companion to the inclusive analysis note\par}
  \vspace{0.5cm}
  {\normalsize\bfseries Secondary analysis under validation --- not part of the inclusive-analysis
    unblinding request\par}
  \vspace{1.5cm}
  {\large Ishanee Pophale and Jarek Nowak\par}
  \vspace{0.5cm}
  {\normalsize\itshape Lancaster University\par}
  \vspace{2cm}
  \begin{abstract}
    The inclusive CC$1\mu1\pi Xp$ measurement of the companion note is repeated on the subsample in
    which a proton candidate with $p_p>0.3\GeVc$ is also reconstructed. Requiring the proton gives
    access to the kinematics of the visible hadronic system: the transverse kinematic imbalance
    between the muon and the $\pi p$ system, a calorimetric hadronic-mass proxy, the proton polar
    angle and the pion--proton opening angle. All of these are extracted in the three NuMI
    configurations with the same Wiener-SVD machinery and released with their additional-smearing
    matrices and covariances. \textbf{Status.} The one-bin proton-tagged total is proposed as a
    secondary result. Every differential and two-dimensional product is a validation product, not
    proposed for approval, until the dedicated inverted-proton-PID control region exists, the
    hadronic-mass proxy and $p_n$ prescriptions are validated for pion production, statistical
    coverage is demonstrated and a non-nominal-model closure is run for this family. The
    differential $W_{\pi p}$ shape is withdrawn. All data points are blind fake data, and the FHC
    and combined normalisations predate the Run-1 exposure correction.
  \end{abstract}
  \vfill
  \tableofcontents
\end{titlepage}

\begin{center}\fbox{\parbox{0.94\textwidth}{\textbf{Product status.}
\textbf{Proposed secondary result:} the proton-tagged one-bin total.
\textbf{Validation only:} $W_{\rm had}$, $\delta p_T$, $\delta\alpha_T$, $\delta\phi_T$, $p_n$,
the inherited muon and pion observables, $\theta_p$, $\theta_{\pi p}$, and the two-dimensional
pairs. \textbf{Withdrawn:} the differential $W_{\pi p}$ shape.
\textbf{Shared with the inclusive note} (not repeated here): exposure, flux, target count,
event weights, detector variations, the unfolding implementation and the covariance conventions.}}
\end{center}

\section{Why a separate measurement}
Requiring a proton changes the measurement, not just the sample. The signal definition acquires a
proton threshold and a leading-proton convention; the reconstruction acquires a proton candidate
that is the right particle $82\%$ of the time; the observables acquire assumptions --- a
single-nucleon residual-mass prescription in the hadronic-mass proxy and in $p_n$ --- that were
written for quasi-elastic kinematics and have not been validated for pion production; and the
detector uncertainty, evaluated on a sample a third the size, reaches $100\%$ in sparse bins. The
inclusive measurement carries none of these, which is why the two are documented and approved
separately.

\section{Proton-tagged selection}
\label{sec:pt_selection}
\label{sec:selection_1p}
\label{sec:wtki_cutflow}

A second selection, the \emph{proton-tagged} one, is applied to the same events. It is
used for the secondary-scope results of \S\ref{sec:wtki}: with an identified proton the
full visible hadronic system is reconstructed, which opens the hadronic invariant mass
and the transverse kinematic imbalance between the muon and the hadronic system. It is
defined here, next to the inclusive selection it extends, so that the two can be read
together.

\paragraph{Definition.} The proton-tagged selection applies the ten inclusive cuts of
\noteref{Sec.}{sec:cutdefs} unchanged and adds an eleventh, the leading-proton identification
(``IdentProton''), which defines the subsample:
\begin{center}\small
\begin{tabular}{lcl}
\toprule
Requirement & Threshold & Description \\
\midrule
Candidate set & --- & track-like primaries other than the muon and pion candidates \\
Proton likeness & LLR PID $<0.05$ & the most proton-like of the set is the candidate \\
Containment & end point in FV & range-based momentum is meaningful only if contained \\
Momentum & $p_p\ge0.3\GeVc$ & tracking threshold; below it protons are not reconstructed reliably \\
\bottomrule
\end{tabular}
\end{center}
\noindent The proton momentum is taken from the proton-hypothesis calorimetric kinetic
energy, $p_p=\sqrt{T_p^2+2m_pT_p}$; truth-matched on the FHC overlay it carries a $-9\%$
bias (underestimate) and a $21\%$ resolution, both folded into the response matrix and
corrected by the unfolding. The signal definition is correspondingly narrowed to the
inclusive signal of \noteref{Sec.}{sec:signal} plus at least one true proton above $0.3\GeVc$;
the muon and pion phase space is unchanged.

\paragraph{Performance.} At full exposure the proton tag yields a selection efficiency
of $11.4\%$ and a purity of $52\%$ in FHC ($50\%$ RHC, $51\%$ combined), against
$17.5\%$ and $\sim59\%$ for the inclusive selection. Relative to the inclusive final cut
evaluated under the proton-tagged signal definition, the tag rejects $67\%$ of the
neutrino background for a $32\%$ signal loss, lifting the purity from $34\%$ to $52\%$
and improving $S/\sqrt{B}$ by $\sim\!20\%$ ($17.1\to20.5$, FHC). Truth-matched, $87\%$
of tagged events contain a true proton above threshold; the $\sim\!13\%$ mis-tag is
dominated by a charged pion or other track taken for the proton, which is why the
proton-tagged results need a control region of their own before they are promoted
(\suppref{Sec.}{sec:cr_1p_scope}). Its clearest benefit is cosmic rejection: requiring a
contained proton above $0.3\GeVc$ removes beam-off background by a factor $3.3$--$3.6$
at the final cut and by up to a factor of six in the CC$0\pi$ control region
(\suppref{Sec.}{sec:sidebands_1p}). The full-exposure cut-flow, stage by stage and per
category, is \suppref{Table}{tab:cutflow_1p} (stacks in \suppref{Fig.}{fig:cutflow_1p});
it uses the same per-run exposure scaling as the inclusive cut-flow of
\noteref{Table}{tab:cutflow}, and its final row is the proton-tagged cut-and-count entry of
\noteref{Table}{tab:cutcount}.


\section{Results, and the status of each}
\label{sec:wtki}

\begin{quote}
\textbf{Scope note.}  This section is secondary scope. Its extractions carry the complete
configured covariance, detector variations included, with detector terms of $10.5\%$ to
$108.5\%$ (median $33.0\%$) against $15.6\%$--$35.0\%$ for the inclusive results: the
same order, as expected for the same detector, except for the six-bin observables on the
sparsest samples, where the unfolding amplifies them. The additional-smearing and
covariance matrices are released with the rest (\noteref{Sec.}{sec:ac_rowsums}). What makes
the scope secondary is a binning limitation that no systematic treatment addresses: the
differential $W_{\pi p}$ shape is withdrawn, because at this exposure no binning
satisfies both the migration and the statistics requirement, so the proton-tagged sample
supports an integrated cross section, not a resonance-shape measurement; and the
proton-mistag background has no dedicated control region yet
(\suppref{Sec.}{sec:cr_1p_scope}). See \noteref{Sec.}{sec:status_matrix}.
\end{quote}

\begin{quote}
\textbf{On the six-bin $W_{\pi p}$ material below.}  Any six-bin $W_{\pi p}$ distribution in this
section is \emph{diagnostic}: it shows why the shape is not measurable and is not a result. The
differential $W_{\pi p}$ shape is withdrawn (\noteref{Appendix}{app:withdrawn}); the proton-tagged
deliverable is the integrated cross section.
\end{quote}
% =============================================================================

Requiring an identified proton in the final state, in addition to the muon and charged
pion, reconstructs the full visible hadronic system and opens two classes of observable
inaccessible to the inclusive $\mathrm{CC}\,1\mu1\pi Xp$ measurement: the \emph{hadronic
invariant mass}, which probes the $\Delta(1232)$ and higher resonances directly, and the
\emph{transverse kinematic imbalance} (TKI) between the muon and the hadronic system,
which is sensitive to Fermi motion and final-state interactions. The proton-tagged
selection, its signal definition and its performance are given in
\S\ref{sec:selection_1p}; in brief, it adds a leading-proton identification to the
inclusive selection, at an efficiency of $11.4\%$ and a purity of $50$--$52\%$, and
trades signal statistics for a fully reconstructed hadronic final state.

\subsection{Extracted cross sections}

Eleven observables are extracted in the proton-tagged subsample; ten are carried as
\emph{secondary validation products}, with their covariances but not proposed for
approval in this review (the binning validation of $W_\mathrm{had}$ and the dedicated
control region are incomplete, \noteref{Sec.}{sec:limitations}); the eleventh, the
differential $W_{\pi p}$ shape, is withdrawn (its files are in the release for
completeness, flagged as such in \texttt{index\_extractions.tsv}), and the robust
proton-tagged deliverable proposed as a headline is the one-bin total cross section
(\noteref{Table}{tab:cutcount}). Two further observables, the proton angles $\theta_p$ and
$\theta_{\pi p}$ (\S\ref{sec:wtki_ang}), were added afterwards and are proposed as secondary scope;
this is the arithmetic behind the twelve rows of Tables~\ref{tab:wtki_fhc}--\ref{tab:wtki_comb},
which are the ten validation products plus those two, the withdrawn $W_{\pi p}$ shape appearing
only in the technical supplement. The ten are: the five hadronic-mass and
transverse-kinematic-imbalance (W/TKI) variables specific to this selection
($W_\mathrm{had}$ and the four TKI observables), plus the five muon/pion kinematic observables of the inclusive analysis
($p_\mu$, $p_\pi$, $\cos\theta_\mu$, $\cos\theta_\pi$, $\theta_{\mu\pi}$) re-measured on
the proton-tagged phase space so that the two selections share a common set of
kinematic distributions. The two hadronic-mass variables are the
\emph{directly measured} $\pi p$ invariant mass,
$W_{\pi p}=\big[(E_\pi+E_p)^2-(\vec p_\pi+\vec p_p)^2\big]^{1/2}$, which makes no
assumption on the neutrino direction, and the \emph{calorimetric} hadronic-mass proxy
$W_\mathrm{had}=\big[m_p^2+2m_p(E_\mathrm{cal}-E_\mu)-Q^2\big]^{1/2}$ built from the
reconstructed calorimetric energy. The four TKI observables; the transverse-momentum
imbalance $\delta p_T$, its two angles $\delta\alpha_T$ and $\delta\phi_T$~\cite{lu_tki}, and the
STV-inferred nucleon-momentum proxy $p_n$~\cite{furmanski_pn} (an approximation, see below), are computed
from the muon versus the summed $(\pi+p)$ hadronic system using the framework STV tool.
Proton momenta are obtained from
the proton-hypothesis range-based kinetic energy. $W_\mathrm{had}$ is measured
in six equal-population bins ($W_{\pi p}$ was binned the same way for the diagnostic
study in the supplement); the four TKI variables use the coarsened two-bin scheme
adopted after the finer binning was found to sit inside the resolution
(\suppref{Sec.}{sec:w_binning}); the reverse-horn-current sample uses the same binning
as FHC. The five inherited kinematic observables use the inclusive-analysis binning (two to seven bins; the proton-tagged
$\cos\theta_\pi$ keeps the former four-bin inclusive edges). All bin edges are set
from the validated true-signal distributions (\suppref{Table}{tab:wtki_binning}).

\paragraph{Reconstruction.} Every observable is built from the three candidate tracks --- muon,
charged pion, proton --- using reconstructed quantities only: range or multiple-scattering
momentum for the muon, range momentum for the pion, calorimetric kinetic energy for the proton,
and the track directions, all in the beam frame. The formulae and the conventions of the
transverse-imbalance tool are given in \suppref{Sec.}{sec:wtki_reco}.

The pipeline is validated end-to-end on FHC with a fake-data closure: the Poisson-thrown
central-value pseudo-data are unfolded with the same Wiener-SVD procedure
(\noteref{Sec.}{sec:unfolding}) and compared to the $A_C$-smeared truth. Table~\ref{tab:wtki_fhc}
shows the result. The bin-by-bin $\chi^2/\mathrm{ndf}$ is well below unity for every
observable, confirming the unfolded pseudo-data are statistically consistent with truth
and that the response, background subtraction and normalisation are self-consistent. The
integrated-cross-section ratios ($0.96$--$1.08$) are ratios to the $A_C$-smeared truth,
so they measure the statistical scatter of this throw (and any bias), not the $A_C$
suppression, which is common to numerator and denominator; the largest departures are
in $\cos\theta_\mu$ and $p_\mu$. The
systematic treatment here includes the flux, cross-section, reinteraction and
detector-variation terms.

The reverse-horn-current measurement is complete at the same binning
(Table~\ref{tab:wtki_rhc}). All ten proton-tagged validation products close positive in RHC, with
$\sigma_\mathrm{int}$ in $0.31$--$0.44\times10^{-38}$~cm$^2$/Ar against $0.39$--$0.50$
for FHC; these are $A_C$-suppressed integrals with different $A_C$ per configuration, so
their ratio is not the RHC/FHC ratio (that is $0.860$ from the one-bin totals,
\noteref{Table}{tab:cutcount}). The closure gives $\chi^2/\mathrm{ndf}\lesssim1$ for almost
every observable.

The two families of processed files store the simulated exposure under different conventions;
the run-by-run normalisation used here is correct under both (\suppref{Sec.}{sec:wtki_reco}).

The combined-sample measurement is complete on the same binning, is included in the
results tables and the data release, and covers all ten combined proton-tagged validation products.

\begin{table}[H]
  \centering
  \caption{FHC fake-data closure for the twelve proton-tagged observables.
    $\sigma_\mathrm{int}$ is the integrated (unfolded) cross section over the
    measured phase space ($10^{-38}$~cm$^2$/Ar); ``unf./truth'' is the ratio of integrals;
    $\chi^2$ is the bin-by-bin comparison of unfolded pseudo-data to $A_C$-smeared truth
    (ndf = number of bins). The upper block is the five released W/TKI observables;
    $W_\mathrm{had}$ is at six bins, while the four TKI variables use the coarsened two-bin
    scheme adopted after the finer binning was found to sit inside the resolution; the
    six-bin $W_{\pi p}$ closure rows are diagnostic only and are given in the technical
    supplement (\suppref{Sec.}{sec:w_binning});
    the lower block is the five kinematic observables inherited from the inclusive
    analysis. The small $\chi^2/\mathrm{ndf}$ validates the chain. The integral ratios compare
    the unfolded result with the $A_C$-smeared truth, so departures from unity reflect
    the statistical realisation of the throw and any residual extraction bias; they do
    not measure the suppression introduced by $A_C$. The third block ($^{\dagger}$) is the two angular observables proposed as secondary scope (\S\ref{sec:wtki_ang}).}
  \label{tab:wtki_fhc}
  \begin{tabular}{lccc}
\toprule
Observable & $\sigma_\mathrm{int}$ & unf./truth & $\chi^2/\mathrm{ndf}$ \\
\midrule
      $W_\mathrm{had}$   & $0.389$ & $1.00$ & $0.88/6$ \\
      $\delta p_T$       & $0.461$ & $0.98$ & $0.01/2$ \\
      $\delta\alpha_T$   & $0.485$ & $0.96$ & $0.02/2$ \\
      $\delta\phi_T$     & $0.474$ & $0.99$ & $0.11/2$ \\
      $p_n$              & $0.490$ & $0.98$ & $0.06/2$ \\
      \midrule
      $p_\mu$            & $0.496$ & $1.04$ & $1.89/7$ \\
      $p_\pi$            & $0.498$ & $0.98$ & $0.01/2$ \\
      $\cos\theta_\mu$   & $0.445$ & $1.08$ & $0.98/5$ \\
      $\cos\theta_\pi$   & $0.394$ & $0.99$ & $0.30/4$ \\
      $\theta_{\mu\pi}$  & $0.469$ & $1.02$ & $0.50/5$ \\
      \midrule
      $\theta_p$$^{\dagger}$         & $0.532$ & $0.98$ & $1.92/4$ \\
      $\theta_{\pi p}$$^{\dagger}$      & $0.456$ & $0.99$ & $0.16/3$ \\
\bottomrule
\end{tabular}
\end{table}


Figure~\ref{fig:dsigma_ccpi1p_fhc} shows the five released FHC W/TKI differential cross
sections ($W_\mathrm{had}$ and the four TKI observables), overlaid with all four generator
predictions; the withdrawn six-bin $W_{\pi p}$ distribution is deliberately not shown
here, and its diagnostic closure and generator curves are in the technical supplement
(\suppref{Sec.}{sec:w_binning}). The generator
W/TKI predictions were regenerated at this binning (GENIE and GiBUU natively, NuWro and
NEUT from their event vectors inside the SL7 container). All curves, data, tune and
generators, are $A_C$-smeared, so the comparison is made consistently in the
measurement space (\suppref{Sec.}{sec:acwtki}). The predictions are built at truth level from each
generator's event-level final state, GENIE gst events, GiBUU final-event records,
and the NuWro and NEUT event vectors reprocessed in their respective SL7-container
environments, by applying the inclusive signal plus a leading true proton above
$0.3$~GeV/$c$ and evaluating the six observables with the same formulas as the selection
(\suppref{Sec.}{sec:wtki_reco}), then combining the $\nu_\mu$ and $\bar\nu_\mu$ samples over the
FHC flux to a per-argon cross section. The four generators span a factor of $\sim1.6$ in
the integrated proton-tagged cross section, GENIE lowest ($0.50$), then NuWro
($0.65$), GiBUU ($0.74$) and NEUT highest ($0.78\times10^{-38}$~cm$^2$/Ar), bracketing
the injected tune (proton-tagged one-bin total, \noteref{Table}{tab:cutcount}). In
the released two-bin TKI scheme the differences between generators are almost entirely
normalisation: their shapes, and that of the tune, agree with the unfolded points within
the uncertainties in every panel, so the present exposure does not discriminate between
the generators' treatments of Fermi motion and final-state interactions. Data and
predictions are both computed about the beam axis (\noteref{Sec.}{sec:beamdir}); a comparison in
which one side used the detector $z$ axis would show apparent shape disagreements of a
factor of several, which is why the frame is stated with every transverse quantity.

\begin{table}[H]
  \centering
  \caption{As Table~\ref{tab:wtki_fhc}, for RHC.}
  \label{tab:wtki_rhc}
  \begin{tabular}{lccc}
\toprule
Observable & $\sigma_\mathrm{int}$ & unf./truth & $\chi^2/\mathrm{ndf}$ \\
\midrule
      $W_\mathrm{had}$   & $0.314$ & $0.96$ & $0.69/6$ \\
      $\delta p_T$       & $0.418$ & $1.10$ & $0.52/2$ \\
      $\delta\alpha_T$   & $0.438$ & $1.08$ & $0.19/2$ \\
      $\delta\phi_T$     & $0.389$ & $1.02$ & $0.10/2$ \\
      $p_n$              & $0.391$ & $1.00$ & $1.06/2$ \\
      \midrule
      $p_\mu$            & $0.380$ & $1.09$ & $3.80/7$ \\
      $p_\pi$            & $0.429$ & $1.09$ & $1.21/2$ \\
      $\cos\theta_\mu$   & $0.392$ & $1.11$ & $0.84/5$ \\
      $\cos\theta_\pi$   & $0.365$ & $1.10$ & $6.01/4$ \\
      $\theta_{\mu\pi}$  & $0.367$ & $0.96$ & $0.14/5$ \\
      \midrule
      $\theta_p$$^{\dagger}$         & $0.391$ & $0.99$ & $0.04/4$ \\
      $\theta_{\pi p}$$^{\dagger}$      & $0.378$ & $1.02$ & $0.12/3$ \\
\bottomrule
\end{tabular}
\end{table}

\begin{table}[H]
  \centering
  \caption{As Table~\ref{tab:wtki_fhc}, for the combined configuration. The combined result is an independent extraction with its own response and $A_C$, not an average of FHC and RHC, so it need not lie between them.}
  \label{tab:wtki_comb}
  \begin{tabular}{lccc}
\toprule
Observable & $\sigma_\mathrm{int}$ & unf./truth & $\chi^2/\mathrm{ndf}$ \\
\midrule
      $W_\mathrm{had}$   & $0.345$ & $0.99$ & $0.85/6$ \\
      $\delta p_T$       & $0.540$ & $1.09$ & $0.17/2$ \\
      $\delta\alpha_T$   & $0.459$ & $1.01$ & $0.02/2$ \\
      $\delta\phi_T$     & $0.503$ & $1.09$ & $0.24/2$ \\
      $p_n$              & $0.449$ & $1.02$ & $0.85/2$ \\
      \midrule
      $p_\mu$            & $0.437$ & $1.03$ & $1.53/7$ \\
      $p_\pi$            & $0.462$ & $1.03$ & $0.40/2$ \\
      $\cos\theta_\mu$   & $0.410$ & $1.06$ & $1.52/5$ \\
      $\cos\theta_\pi$   & $0.382$ & $1.02$ & $3.40/4$ \\
      $\theta_{\mu\pi}$  & $0.445$ & $1.04$ & $0.26/5$ \\
      \midrule
      $\theta_p$$^{\dagger}$         & $0.485$ & $1.01$ & $0.67/4$ \\
      $\theta_{\pi p}$$^{\dagger}$      & $0.413$ & $1.01$ & $0.06/3$ \\
\bottomrule
\end{tabular}
\end{table}

\begin{figure}[H]
  \centering
  \blindfig[\textwidth]{dsigma_ccpi1p_noW_FHC5}
  \caption{Proton-tagged released differential cross sections for
    FHC: the calorimetric hadronic mass
    $W_\mathrm{had}$ and the four transverse-kinematic-imbalance observables
    $\delta p_T$, $\delta\alpha_T$, $\delta\phi_T$ and $p_n$ (the withdrawn $W_{\pi p}$
    shape is not shown; see the technical supplement). Unfolded fake data
    (black points) are compared to the $A_C$-smeared fake-data truth (grey dashed), the
    MicroBooNE tune (black), and the GENIE (blue), GiBUU (green), NEUT (purple) and
    NuWro (orange) predictions built for the same proton-tagged phase space.
    $\delta\alpha_T$ rises toward $180^\circ$, the
    signature of final-state interactions.}
  \label{fig:dsigma_ccpi1p_fhc}
\end{figure}

\begin{figure}[H]
  \centering
  \blindfig[\textwidth]{dsigma_ccpi1p_noW_RHCFULL}
  \caption{The same five proton-tagged observables for RHC. Curves and points as in
    Fig.~\ref{fig:dsigma_ccpi1p_fhc}.}
  \label{fig:dsigma_ccpi1p_rhc}
\end{figure}

\begin{figure}[H]
  \centering
  \blindfig[\textwidth]{dsigma_ccpi1p_noW_COMB}
  \caption{The same five proton-tagged observables for the combined configuration. Curves and
    points as in Fig.~\ref{fig:dsigma_ccpi1p_fhc}.}
  \label{fig:dsigma_ccpi1p_comb}
\end{figure}


With all three configurations in hand, the combined $\sigma_\mathrm{int}$ lies between the
FHC and RHC values for eight of the ten proton-tagged validation products; $\delta p_T$ ($0.540$
against $0.461$ and $0.418$) and $\delta\phi_T$ ($0.503$ against $0.474$ and $0.389$)
lie above both, by amounts well inside their $\sim\!30\%$ uncertainties.

That is not a defect: the combined measurement is \emph{not} required to lie between the two
single-mode ones.  It is an independent extraction with its own response matrix, its own
regularisation and its own $A_C$, not a fluence-weighted average, so excursions at the
few-per-cent level would have been unremarkable against a $\sim\!30\%$ systematic
uncertainty.

The spread across observables within a configuration, $0.35$ to $0.54$ for the combined
set, is likewise not a physics spread but the residual of the additional smearing
discussed in the next section: each observable's $A_C$ rescales the integral by a different
factor, so the ten numbers are not ten independent measurements of the same quantity.

\subsection{What the proton candidate is}
\label{sec:prmatch}

Every proton-tagged observable is built on the reconstructed proton candidate, while the truth-level
definitions use the leading true proton. On selected signal events the candidate is the leading true
proton $81.7\%$ of the time, some true proton $87.4\%$, and not a proton at all $12.6\%$
($10\,105$ events).

What that costs was measured for the two angular observables, and the answer depends on which
question is asked. A correct pick reconstructs the angle tightly (interquartile range $0.073$~rad in
$\theta_p$, $0.124$~rad in $\theta_{\pi p}$); a wrong pick does not smear the angle, it randomises it
($1.0$--$1.4$~rad, three quarters of those events beyond $0.3$~rad). Averaged over everything the
candidate widens the \emph{core} resolution by only a quarter to a third, which understates it badly,
because a binned measurement is not limited by its core. Residuals beyond $0.3$~rad occur in $18.9\%$
of $\theta_p$ events and \emph{$74\%$ of those are wrong picks} ($\theta_{\pi p}$: $25.9\%$ and
$51\%$). Candidate purity is therefore the dominant source of bin-to-bin migration in $\theta_p$,
which is why its migration diagonal sits at the criterion and its RHC response is rank-deficient.

The same $18\%$ wrong-pick rate applies to $W_{\pi p}$, $W_\mathrm{had}$ and the TKI observables,
which use the same candidate; the migration attribution above was measured for the angles only. The
control region that would constrain this is the multi-$\pi$ sideband, which is already open:
\texttt{pion\_number} is defined by the same PID cut that picks the proton, so an otherwise-signal
event whose proton fails it has two pion candidates and lands there. It constrains the
proton~$\to$~pion direction rather than the pion~$\to$~proton one that produces the $12.6\%$, and it
is the most signal-contaminated of the four regions, so a region inverting the proton PID at exactly
one pion candidate is what would separate them.

\subsection{Proton angular observables: \texorpdfstring{$\theta_p$}{theta p} and
  \texorpdfstring{$\theta_{\pi p}$}{theta pi p}}
\label{sec:wtki_ang}

Two angular observables of the proton-tagged sample are proposed as secondary scope. Like every
other number in this note they are fake data. Both use the leading proton candidate:
\[
  \theta_p=\arccos\!\big(\hat p_p\cdot\hat\nu\big), \qquad
  \theta_{\pi p}=\arccos\!\big(\hat p_\pi\cdot\hat p_p\big),
\]
with $\hat\nu$ the beam-frame neutrino direction of \noteref{Sec.}{sec:beamdir} --- the same convention as
every other angle here --- while $\theta_{\pi p}$ is an angle between two reconstructed tracks and
so needs no beam direction at all, which makes it the one angular observable in this analysis that
is insensitive to the beam-frame convention. $\theta_p$ is available from the stored proton
direction; $\theta_{\pi p}$ required the proton-tagged tree to be reprocessed to write
\texttt{pi\_pr\_opening\_angle\_\{reco,true\}}, checked to leave every existing branch unchanged.

They measure something the released set does not. The five W/TKI observables are built from the
muon against the \emph{summed} hadronic system, so they are insensitive to how that system shares
its momentum between pion and proton; $\theta_{\pi p}$ is exactly that sharing. $\theta_p$ is the
proton direction alone, where the released set has only the derived quantities.

Both unfold well (Tables~\ref{tab:wtki_fhc}--\ref{tab:wtki_comb}). $\theta_{\pi p}$ is the
best-conditioned observable of the proton-tagged programme, $s_{\min}/s_1=0.30$--$0.48$ against
$0.27/0.24$ for the released $\cos\theta_\pi$, at the price of a last bin whose total uncertainty
reaches $93\%$ (FHC) and $187\%$ (RHC). $\theta_p$ conditions at $0.383$ (FHC) and $0.408$
(combined) but is rank-deficient in RHC ($0.001$), so it is proposed for FHC and the combined
configuration only; coarsening the RHC binning was tried and does not rescue it. Its edges are
derived on the RHC chain rather than the FHC one: at the same four bins they condition three times
better in both usable configurations. Both distributions are strongly forward peaked, and the
generators differ mainly in normalisation rather than shape, so these observables discriminate
through the integral. The study behind all of this --- binning, generator predictions, the
$0.50$-criterion screen and the two-dimensional programme --- is a separate note.

\begin{figure}[H]
  \centering
  \blindfig[\textwidth]{newobs_xsec}
  \caption{The two proton angular observables in all three configurations: $\theta_p$ (top row) and
    $\theta_{\pi p}$ (bottom row), for FHC, RHC and the combined exposure. Points are the unfolded
    fake data with the total uncertainty, the grey histogram is the $A_C$-smeared fake-data truth,
    and the four coloured histograms are the generator predictions; integrated cross sections are
    given in the legends. The RHC $\theta_p$ panel is shown for completeness only: its response is
    rank-deficient and its shape is not a measurement.}
  \label{fig:newobs_xsec}
\end{figure}

\subsection{Two-dimensional cross sections (proposed)}
\label{sec:2d}

Eight two-dimensional pairs were screened and extracted in all three configurations. The framework
needed no new unfolding machinery --- a two-dimensional binning is a flat list of analysis bins
ordered slice by slice --- and the integrated normalisation of the two-dimensional chain agrees with
the released one-bin total to a median of $0.7\%$ once each is compared with its own $A_C$-smeared
truth. Two pairs are worth proposing, both in the combined configuration only:
$\theta_p\times\delta p_T$ ($s_{\min}/s_1=0.617$) and $\cos\theta_\pi\times\cos\theta_\mu$
($0.563$). Combining the beam modes roughly doubles their conditioning ($0.46\to0.62$,
$0.25\to0.56$), which is the clearest benefit the joint exposure delivers anywhere in this analysis,
and it comes from statistics and the shared flux rather than from any $\nu/\bar\nu$ separation.
Everything built on $W_{\pi p}$ is unusable, as its $39\%$ leakage between outer bins predicted.

One pair is worth recording as a caution. $\theta_{\pi p}\times W_{\pi p}$ conditions at $0.560$,
as well as the two usable pairs, and its extraction is still worthless: a $58\%$ data-statistical
term, per-bin uncertainties to $545\%$ and a closure ratio of $-0.40$. Conditioning bounds how much
of the response can be inverted and says nothing about how much data survives in the inverted
directions; the two have to be read together, here and everywhere else in this note.

\begin{table}[H]
  \centering
  \caption{The eight two-dimensional pairs (inner $\times$ outer), ordered by the conditioning of
    $A_C$ in the combined configuration. The last three columns are the combined configuration
    only: the bin-averaged data-statistical term, the range of the per-bin total uncertainty and the
    range of the closure ratio against the $A_C$-smeared truth. Fake data.}
  \label{tab:2dextract}
  \small
  \begin{tabular}{lcccccc}
    \toprule
    & \multicolumn{3}{c}{$s_{\min}/s_1$} & \multicolumn{3}{c}{combined configuration} \\
    \cmidrule(lr){2-4}\cmidrule(lr){5-7}
    Pair & FHC & RHC & COMB & data stat. & unc.\ [\%] & closure \\
    \midrule
    $\theta_p\times\delta p_T$ & $0.455$ & $0.332$ & $\mathbf{0.617}$ & $13.7\%$ & $31$--$88$ & $0.82$--$1.25$ \\
    $\cos\theta_\pi\times\cos\theta_\mu$ & $0.245$ & $0.315$ & $\mathbf{0.563}$ & $14.8\%$ & $32$--$62$ & $0.93$--$1.17$ \\
    $\theta_{\pi p}\times W_{\pi p}$ & $0.403$ & $0.214$ & $0.560$ & $57.7\%$ & $43$--$545$ & $-0.40$--$1.10$ \\
    $\cos\theta_\pi\times\theta_{\mu\pi}$ & $0.027$ & $0.133$ & $0.169$ & $22.2\%$ & $28$--$94$ & $0.91$--$1.37$ \\
    $\cos\theta_\mu\times p_\mu$ & $0.002$ & $0.261$ & $0.110$ & $8.9\%$ & $33$--$57$ & $0.97$--$0.99$ \\
    $\theta_p\times W_{\pi p}$ & $0.008$ & $0.152$ & $0.105$ & $111.7\%$ & $36$--$291$ & $0.39$--$1.22$ \\
    $\cos\theta_\pi\times p_\pi$ & $0.023$ & $0.158$ & $0.052$ & $12.9\%$ & $32$--$58$ & $0.96$--$1.14$ \\
    $\theta_{\pi p}\times\delta p_T$ & $0.024$ & $0.035$ & $0.028$ & $10.7\%$ & $28$--$50$ & $0.93$--$1.08$ \\
    \bottomrule
  \end{tabular}
\end{table}

Two entries in that table are worth reading against each other. $\cos\theta_\mu\times p_\mu$
screened best of the eight on the migration diagonal ($0.70/0.68$) and is rank-deficient in FHC
($0.002$); $\theta_p\times\delta p_T$ screened worst of the proton-tagged pairs and conditions best
of all. The diagonal is a necessary condition, not a sufficient one, and the regularised response is
what decides --- the same lesson the one-dimensional binning study returns
(\noteref{Appendix}{app:altbinning}).

Figures~\ref{fig:2dslices} and~\ref{fig:2dmap} show the two usable pairs. The slice form is how a
double-differential result is normally published and how the generators are compared; the map is the
same numbers, and shows at a glance where the measurement has support. Each slice panel carries its
own vertical scale, because the outer bins differ by an order of magnitude, and the maps are drawn
on equal-size cells labelled by their ranges for the same reason.

\begin{figure}[H]
  \centering
  \blindfig[0.66\textwidth]{bkgfit/xsec2d_thetap_dpt_COMB_slices}\\[6pt]
  \blindfig[0.98\textwidth]{bkgfit/xsec2d_costhpi_costhmu_COMB_slices}
  \caption{The two proposed two-dimensional measurements as slices in the outer variable, combined
    configuration: $\theta_p\times\delta p_T$ (top) and $\cos\theta_\pi\times\cos\theta_\mu$
    (bottom). Points are the unfolded fake data with the total uncertainty, the grey histogram is
    the $A_C$-smeared fake-data truth and the four coloured histograms are the generator
    predictions.}
  \label{fig:2dslices}
\end{figure}

\begin{figure}[H]
  \centering
  \blindfig[0.82\textwidth]{bkgfit/xsec2d_thetap_dpt_COMB_map}\\[6pt]
  \blindfig[0.90\textwidth]{bkgfit/xsec2d_costhpi_costhmu_COMB_map}
  \caption{The same two measurements as maps: $\mathrm{d}^2\sigma/\mathrm{d}X\,\mathrm{d}Y$ per cell
    with its total uncertainty (left) and the ratio to the $A_C$-smeared truth (right), for
    $\theta_p\times\delta p_T$ (top) and $\cos\theta_\pi\times\cos\theta_\mu$ (bottom). Cells are
    equal-size and labelled by their ranges; the colour scale is the one used for the
    matrix figures elsewhere in this note. The closure panels are flat to within their
    uncertainties, which are the measurement uncertainty divided by the truth.}
  \label{fig:2dmap}
\end{figure}

The per-mode extractions of these two pairs are not shape measurements --- their responses cannot be
inverted --- and are not shown. Whether a combined-only result is acceptable, when every released
differential result is also quoted per beam mode, is the open question these two pairs raise.

% =============================================================================

\section{Limitations and approval prerequisites}
\label{sec:pt_prereq}
Before any proton-tagged differential result is promoted beyond validation:
\begin{enumerate}\itemsep2pt
  \item construct and open the inverted-proton-PID control region, which separates the
    proton$\to$pion and pion$\to$proton misidentification directions;
  \item validate the $W_{\rm had}$ binning with an edge scan, and the pion-production meaning of
    $p_n$, $p_L$ and the calorimetric mass proxy;
  \item demonstrate statistical coverage with ensembles for this family;
  \item run a non-nominal-model closure for this family;
  \item bound the candidate-mistag effect, which supplies $74\%$ of the large migrations in
    $\theta_p$;
  \item reassess whether any two-dimensional result has the effective rank and statistical
    support to be a measurement.
\end{enumerate}

% glossary.tex -- shared by the analysis note and the technical supplement.
% Plain-language definitions, in the order a reader meets the ideas, not alphabetical.
\section*{Glossary}
\addcontentsline{toc}{section}{Glossary}
\label{sec:glossary}

\begin{description}\itemsep2pt
  \item[Signal definition] The set of particles leaving the nucleus that counts as the process being
    measured, with momentum thresholds and a fiducial requirement. It is deliberately \emph{not} a
    statement about the interaction mechanism, because final-state interactions mix mechanisms and
    only the final state is observable.
  \item[Fake data / pseudo-data] A simulated dataset used in place of beam-on data while the
    analysis is blind. Here it is a Poisson throw: each simulated event is replicated a
    Poisson-distributed number of times with mean equal to its weight at data exposure, so the
    sample has the statistical fluctuations of a real dataset.
  \item[Blind] The signal region has not been looked at in beam-on data. Every plotted ``data''
    point in these documents is fake data.
  \item[Beam-on, beam-off (EXT), dirt] Beam-on data are recorded in the beam spill; beam-off data are
    recorded between spills and measure cosmic-ray activity; ``dirt'' is simulated neutrino
    interactions outside the cryostat whose products enter the detector. The prediction is
    neutrino simulation $+$ beam-off $+$ dirt.
  \item[Gate ratio] The factor scaling beam-off data to the beam-on exposure: the ratio of the number
    of beam-on to beam-off trigger gates in a running period.
  \item[Tune, central value (CV)] The \uboone{} tune is the collaboration's adjusted GENIE model; the
    central value is the nominal prediction from it. ``CV weight'' is the per-event weight that
    turns the raw simulation into the tuned prediction.
  \item[Safe-weight rule] Any event weight that is non-finite, negative or above $30$ is replaced by
    $1$. One rule, applied everywhere, so that the prediction and every pseudo-data throw agree.
  \item[Universe] One alternative version of the prediction with a systematic parameter varied. A
    \emph{multisim} is a set of hundreds of universes drawn randomly from the parameter
    uncertainties; a \emph{unisim} is a single $\pm1\sigma$ variation of one parameter. PPFX
    universes vary the flux; GENIE universes the interaction model; reinteraction universes the
    hadron transport.
  \item[Detector variations (detVar)] Dedicated simulation samples with one aspect of the detector
    response changed --- light yield, recombination, space-charge, wire response --- used to
    estimate the detector uncertainty.
  \item[Migration matrix, efficiency, smearceptance] The migration matrix $M_{ij}$ is the probability
    that an event in true bin $j$ is reconstructed in bin $i$; the efficiency $\varepsilon_j$ is the
    fraction of true bin $j$ that is selected at all. Their product $S_{ij}=\varepsilon_j M_{ij}$
    is the response the unfolding inverts; we call it the smearceptance matrix.
  \item[Migration diagonal] The fraction of events reconstructed in a bin that were truly in that
    bin (the column-normalised diagonal of $M$). A binning is accepted only if every bin's diagonal
    exceeds $0.68$.
  \item[Unfolding, Wiener-SVD] Solving for the true spectrum given the reconstructed one and the
    response. A plain inversion amplifies statistical noise; Wiener-SVD regularises it with a filter
    built from the simulation, trading a small bias for a large reduction in variance.
  \item[Additional-smearing matrix $A_C$] The price of regularisation, made explicit: the unfolded
    result estimates $A_C\,T$ rather than the truth $T$. A prediction must be multiplied by $A_C$
    before it is compared with the result. $A_C$ is released with every extraction.
  \item[Not normalisation-preserving] The rows of $A_C$ do not sum to one, so the integral of an
    unfolded differential result is not the total cross section and differs between observables of
    the same sample. The total is measured separately.
  \item[One-bin extraction] The total cross section, extracted with a single true bin covering the
    whole phase space. Nothing is regularised, $A_C=1$, and the result is the physical total.
  \item[Conditioning, $s_{\min}/s_1$, rank-deficient] The ratio of the smallest to the largest
    singular value of the regularised response. Near zero, some direction in the true spectrum
    cannot be recovered from the data at all (rank-deficient); the result in that direction is the
    prior. A binning can pass the migration-diagonal criterion and still be unusable here.
  \item[Closure] Unfolding pseudo-data and checking that the truth comes back. Closure against the
    $A_C$-smeared truth tests the implementation; it cannot test whether the model is right,
    because the same model is on both sides.
  \item[Covariance] The matrix of uncertainties and their correlations between bins, one per
    systematic source, propagated through the unfolding. ``Nominal configured covariance'' means the
    sum over the sources implemented in the analysis; the sources known to be missing are listed
    in the limitations.
  \item[Flux-averaged cross section] The cross section averaged over the neutrino energy spectrum
    of the beam, per argon nucleus, in the stated phase space. It depends on the flux, which is why
    each horn configuration is a separate measurement.
  \item[FHC, RHC] Forward and reverse horn current: the two polarities of the NuMI focusing horns,
    giving a neutrino-dominated and an antineutrino-enriched beam.
  \item[Beam frame] All angles are measured about the neutrino direction, which at \uboone{}'s
    off-axis position is $28^\circ$ from the detector $z$ axis.
  \item[Control region, sideband] A sample selected by inverting one requirement of the signal
    selection, so that it is disjoint from the signal region and enriched in one background. It
    can be inspected on beam-on data while the signal region stays blind.
  \item[Transverse kinematic imbalance (TKI)] Variables built from the momenta transverse to the
    neutrino direction of the lepton and the hadronic system: $\delta p_T$, $\delta\alpha_T$,
    $\delta\phi_T$ and the inferred nucleon momentum $p_n$. They isolate nuclear effects from the
    unknown neutrino energy.
  \item[LLR particle identification] The log-likelihood-ratio score built from track $dE/dx$ that
    separates protons from muons and pions.
  \item[Generator predictions (FTE files)] Standalone GENIE, NuWro, NEUT and GiBUU predictions,
    run with the NuMI flux and binned identically to the results, stored as histograms for the
    comparison figures.
\end{description}

\begin{thebibliography}{99}

\bibitem{numi_nue_ccpi}
  MicroBooNE Collaboration,
  ``Measurement of the flux-averaged differential $\nu_e$ CC$1\pi$ cross section
  in the NuMI beam,'' Internal Note v3.2 (2025).

\bibitem{mistry_thesis}
  K.~Mistry,
  ``First Measurement of the Flux-Averaged Differential Charged-Current
  Electron-Neutrino and Antineutrino Cross Section on Argon with the MicroBooNE
  Detector,'' PhD thesis (2021).

\bibitem{nova}
  NOvA Collaboration, M.~A.~Acero \textit{et al.},
  ``New constraints on oscillation parameters from $\nu_e$ appearance and
  $\nu_\mu$ disappearance in the NOvA experiment,''
  \textit{Phys.\ Rev.\ D} \textbf{98}, 032012 (2018).

\bibitem{t2k}
  T2K Collaboration, K.~Abe \textit{et al.},
  ``Constraint on the matter--antimatter symmetry-violating phase in neutrino
  oscillations,''
  \textit{Nature} \textbf{580}, 339--344 (2020).

\bibitem{dune}
  DUNE Collaboration, B.~Abi \textit{et al.},
  ``Deep Underground Neutrino Experiment (DUNE), Far Detector Technical Design
  Report, Volume II: DUNE Physics,'' arXiv:2002.03005 [hep-ex] (2020).

\bibitem{pandora}
  J.~S.~Marshall and M.~A.~Thomson,
  ``The Pandora Software Development Kit for Pattern Recognition,''
  \textit{Eur.\ Phys.\ J.\ C} \textbf{75}, 439 (2015).

\bibitem{wienervsd}
  W.~Tang, X.~Li, X.~Qian, H.~Wei, and C.~Zhang,
  ``Data Unfolding with Wiener-SVD Method,''
  \textit{JINST} \textbf{12}, P10002 (2017), arXiv:1705.03568 [physics.data-an].

\bibitem{ppfx}
  L.~Aliaga \textit{et al.}\ (MINERvA Collaboration),
  ``Neutrino Flux Predictions for the NuMI Beam,''
  \textit{Phys.\ Rev.\ D} \textbf{94}, 092005 (2016).

\bibitem{genie}
  C.~Andreopoulos \textit{et al.},
  ``The GENIE Neutrino Monte Carlo Generator,''
  \textit{Nucl.\ Instrum.\ Methods A} \textbf{614}, 87--104 (2010).

\bibitem{sce}
  MicroBooNE Collaboration, P.~Abratenko \textit{et al.},
  ``Measurement of space charge effects in the MicroBooNE LArTPC using cosmic
  muons,''
  \textit{JINST} \textbf{15}, P12037 (2020).

\bibitem{uboone_laser}
  MicroBooNE Collaboration, C.~Adams \textit{et al.},
  ``A Method to Determine the Electric Field of Liquid Argon Time Projection
  Chambers Using a UV Laser System and Its Application in MicroBooNE,''
  \textit{JINST} \textbf{15}, P07010 (2020).

\bibitem{uboone_pandora}
  MicroBooNE Collaboration, R.~Acciarri \textit{et al.},
  ``The Pandora multi-algorithm approach to automated pattern recognition of
  cosmic-ray muon and neutrino events in the MicroBooNE detector,''
  \textit{Eur.\ Phys.\ J.\ C} \textbf{78}, 82 (2018).

\bibitem{uboone_llr}
  MicroBooNE Collaboration, P.~Abratenko \textit{et al.},
  ``Calorimetric classification of track-like signatures in liquid argon TPCs
  using MicroBooNE data,''
  \textit{JHEP} \textbf{12}, 153 (2021).

\bibitem{uboone_mcs}
  MicroBooNE Collaboration, P.~Abratenko \textit{et al.},
  ``Determination of muon momentum in the MicroBooNE LArTPC using an improved
  model of multiple Coulomb scattering,''
  \textit{JINST} \textbf{12}, P10010 (2017).

\bibitem{uboone_detvar}
  MicroBooNE Collaboration, P.~Abratenko \textit{et al.},
  ``Novel approach for evaluating detector-related uncertainties in a LArTPC
  using MicroBooNE data,''
  \textit{Eur.\ Phys.\ J.\ C} \textbf{82}, 454 (2022).

\bibitem{uboone_tune}
  MicroBooNE Collaboration, P.~Abratenko \textit{et al.},
  ``New CC0$\pi$ GENIE model tune for MicroBooNE,''
  \textit{Phys.\ Rev.\ D} \textbf{105}, 072001 (2022), arXiv:2110.14028 [hep-ex].

\bibitem{uboone_tki}
  MicroBooNE Collaboration, P.~Abratenko \textit{et al.},
  ``First double-differential measurement of kinematic imbalance in neutrino
  interactions with the MicroBooNE detector,''
  \textit{Phys.\ Rev.\ Lett.} \textbf{131}, 101802 (2023), arXiv:2301.03700 [hep-ex].

\bibitem{xsec_analyzer}
  S.~Gardiner \textit{et al.},
  ``XSecAnalyzer: a cross-section analysis framework for MicroBooNE,''
  \url{https://github.com/uboone/xsec_analyzer} (accessed 2026).

\bibitem{numi_beam}
  P.~Adamson \textit{et al.},
  ``The NuMI Neutrino Beam,''
  \textit{Nucl.\ Instrum.\ Methods A} \textbf{806}, 279--306 (2016).

\bibitem{genie3}
  L.~Alvarez-Ruso \textit{et al.}\ (GENIE Collaboration),
  ``Recent highlights from GENIE v3,''
  \textit{Eur.\ Phys.\ J.\ ST} \textbf{230}, 4449--4467 (2021).

\bibitem{nuwro}
  T.~Golan, J.~T.~Sobczyk, and J.~\.Zmuda,
  ``NuWro: the Wroc\l{}aw Monte Carlo generator of neutrino interactions,''
  \textit{Nucl.\ Phys.\ B Proc.\ Suppl.} \textbf{229--232}, 499 (2012).

\bibitem{gibuu}
  O.~Buss \textit{et al.},
  ``Transport-theoretical description of nuclear reactions,''
  \textit{Phys.\ Rept.} \textbf{512}, 1--124 (2012);
  U.~Mosel, ``Neutrino event generators: foundation, status and future,''
  \textit{J.\ Phys.\ G} \textbf{46}, 113001 (2019).

\bibitem{neut2021}
  Y.~Hayato and L.~Pickering,
  ``The NEUT neutrino interaction simulation program library,''
  \textit{Eur.\ Phys.\ J.\ ST} \textbf{230}, 4469--4481 (2021).

\bibitem{geant4}
  S.~Agostinelli \textit{et al.}\ (GEANT4 Collaboration),
  ``GEANT4: a simulation toolkit,''
  \textit{Nucl.\ Instrum.\ Methods A} \textbf{506}, 250--303 (2003).

\bibitem{geant4reweight}
  J.~Calcutt, C.~Thorpe, K.~Mahn, and L.~Fields,
  ``Geant4Reweight: a framework for evaluating and propagating hadronic
  interaction uncertainties in Geant4,''
  \textit{JINST} \textbf{16}, P08042 (2021).

\bibitem{tmva}
  A.~Hoecker \textit{et al.},
  ``TMVA: Toolkit for Multivariate Data Analysis,''
  PoS \textbf{ACAT}, 040 (2007), arXiv:physics/0703039.

\bibitem{dagostini}
  G.~D'Agostini,
  ``A multidimensional unfolding method based on Bayes' theorem,''
  \textit{Nucl.\ Instrum.\ Methods A} \textbf{362}, 487--498 (1995).

\bibitem{lu_tki}
  X.-G.~Lu, L.~Pickering, S.~Dolan, G.~Barr, D.~Coplowe, Y.~Uchida, D.~Wark,
  M.~O.~Wascko, A.~Weber, and T.~Yuan,
  ``Measurement of nuclear effects in neutrino interactions with minimal
  dependence on neutrino energy,''
  \textit{Phys.\ Rev.\ C} \textbf{94}, 015503 (2016).

\bibitem{furmanski_pn}
  A.~P.~Furmanski and J.~T.~Sobczyk,
  ``Neutrino energy reconstruction from one-muon and one-proton events,''
  \textit{Phys.\ Rev.\ C} \textbf{95}, 065501 (2017).

\bibitem{rs_coherent}
  D.~Rein and L.~M.~Sehgal,
  ``Coherent $\pi^0$ production in neutrino reactions,''
  \textit{Nucl.\ Phys.\ B} \textbf{223}, 29--44 (1983).

\bibitem{berger_sehgal_coherent}
  C.~Berger and L.~M.~Sehgal,
  ``Partially conserved axial vector current and coherent pion production by
  low energy neutrinos,''
  \textit{Phys.\ Rev.\ D} \textbf{79}, 053003 (2009).

\bibitem{genieccpi}
  MicroBooNE Collaboration, P.~Abratenko \textit{et al.},
  ``First Measurement of Energy-Dependent Inclusive Muon-Neutrino
  Charged-Current Cross Sections on Argon with the MicroBooNE Detector,''
  \textit{Phys.\ Rev.\ Lett.} \textbf{128}, 151801 (2022), arXiv:2110.14023 [hep-ex].

\bibitem{uboone_nim}
  MicroBooNE Collaboration, R.~Acciarri \textit{et al.},
  ``Design and Construction of the MicroBooNE Detector,''
  \textit{JINST} \textbf{12}, P02017 (2017).

\bibitem{uboone_crt}
  MicroBooNE Collaboration, C.~Adams \textit{et al.},
  ``Design and construction of the MicroBooNE Cosmic Ray Tagger system,''
  \textit{JINST} \textbf{14}, P04004 (2019).

\bibitem{rs1981}
  D.~Rein and L.~M.~Sehgal,
  ``Neutrino Excitation of Baryon Resonances and Single Pion Production,''
  \textit{Ann.\ Phys.} \textbf{133}, 79--153 (1981).

\bibitem{kuzmin2004}
  K.~S.~Kuzmin, V.~V.~Lyubushkin, and V.~A.~Naumov,
  ``Lepton polarization in neutrino nucleon interactions,''
  \textit{Mod.\ Phys.\ Lett.\ A} \textbf{19}, 2815--2829 (2004).

\bibitem{graczyk2008}
  K.~M.~Graczyk and J.~T.~Sobczyk,
  ``Form factors in the quark resonance model,''
  \textit{Phys.\ Rev.\ D} \textbf{77}, 053001 (2008)
  [Erratum: \textit{Phys.\ Rev.\ D} \textbf{79}, 079903 (2009)].

\bibitem{berger2007}
  C.~Berger and L.~Sehgal,
  ``Lepton mass effects in single pion production by neutrinos,''
  \textit{Phys.\ Rev.\ D} \textbf{76}, 113004 (2007).

\bibitem{nowak2009}
  J.~A.~Nowak (MiniBooNE Collaboration),
  ``Four-momentum transfer discrepancy in the charged current $\pi^+$ production
  in the MiniBooNE: Data vs.\ theory,''
  \textit{AIP Conf.\ Proc.} \textbf{1189}, 243--248 (2009).

\bibitem{neut}
  Y.~Hayato,
  ``A neutrino interaction simulation program library NEUT,''
  \textit{Acta Phys.\ Polon.\ B} \textbf{40}, 2477--2489 (2009).

\bibitem{anl_radecky1982}
  G.~M.~Radecky \textit{et al.}\ (ANL Collaboration),
  ``Study of single pion production by weak charged currents in low-energy
  $\nu d$ interactions,''
  \textit{Phys.\ Rev.\ D} \textbf{25}, 1161--1173 (1982).

\bibitem{bnl_kitagaki1986}
  T.~Kitagaki \textit{et al.}\ (BNL Collaboration),
  ``Charged-current exclusive pion production in neutrino--deuterium
  interactions,''
  \textit{Phys.\ Rev.\ D} \textbf{34}, 2554--2565 (1986).

\bibitem{k2k_rodriguez2008}
  A.~Rodriguez \textit{et al.}\ (K2K Collaboration),
  ``Measurement of single charged pion production in the charged-current
  interactions of neutrinos in a 1.3~GeV wide band beam,''
  \textit{Phys.\ Rev.\ D} \textbf{78}, 032003 (2008).

\bibitem{miniboone_ccpi}
  A.~A.~Aguilar-Arevalo \textit{et al.}\ (MiniBooNE Collaboration),
  ``Measurement of $\nu_\mu$-induced charged-current single positive pion
  production cross sections on mineral oil at $E_\nu \in 0.5$--$2.0$~GeV,''
  \textit{Phys.\ Rev.\ D} \textbf{83}, 052007 (2011).

\bibitem{minerva_tki}
  X.-G.~Lu \textit{et al.}\ (MINERvA Collaboration),
  ``Measurement of final-state correlations in neutrino muon-proton mesonless
  production on hydrocarbon at $\langle E_\nu\rangle=3$~GeV,''
  Phys.\ Rev.\ Lett.\ \textbf{121}, 022504 (2018).

\bibitem{minerva_ccpi}
  B.~Eberly \textit{et al.}\ (MINERvA Collaboration),
  ``Charged pion production in $\nu_\mu$ interactions on hydrocarbon at
  $\langle E_\nu \rangle = 4.0$~GeV,''
  \textit{Phys.\ Rev.\ D} \textbf{92}, 092008 (2015).

\bibitem{minerva_fe}
  C.~L.~McGivern \textit{et al.}\ (MINERvA Collaboration),
  ``Cross sections for $\nu_\mu$ and $\bar{\nu}_\mu$ induced pion production on
  hydrocarbon in the few-GeV region using MINERvA,''
  \textit{Phys.\ Rev.\ D} \textbf{94}, 052005 (2016).

\bibitem{minerva_lownu}
  MINERvA Collaboration, J.~Wolcott \textit{et al.},
  ``Evidence for Quasielastic Charge Current Antineutrino Scattering Off
  Protons in Hydrocarbon at $\langle E_{\bar\nu} \rangle = 3.5$~GeV,''
  \textit{Phys.\ Rev.\ Lett.} \textbf{116}, 081802 (2016).

\bibitem{t2k_ccpip}
  K.~Abe \textit{et al.}\ (T2K Collaboration),
  ``Measurement of the muon neutrino charged-current single $\pi^+$ production
  on hydrocarbon using the T2K off-axis near detector ND280,''
  \textit{Phys.\ Rev.\ D} \textbf{101}, 012007 (2020), arXiv:1909.03936 [hep-ex].

\bibitem{argoneut_ccpi}
  R.~Acciarri \textit{et al.}\ (ArgoNEUT Collaboration),
  ``First measurement of the cross section for $\nu_\mu$ and $\bar{\nu}_\mu$
  induced single charged pion production on argon using ArgoNEUT,''
  \textit{Phys.\ Rev.\ D} \textbf{98}, 052002 (2018), arXiv:1804.10294 [hep-ex].

\bibitem{uboone_nue_ccpi}
  MicroBooNE Collaboration, P.~Abratenko \textit{et al.},
  ``First measurement of $\nu_e$ and $\bar{\nu}_e$ charged-current single
  charged-pion production differential cross sections on argon using the
  MicroBooNE detector,''
  \textit{Phys.\ Rev.\ Lett.} \textbf{135}, 061802 (2025),
  arXiv:2503.23384 [hep-ex].

\bibitem{uboone_bnb2026_ccpi}
  MicroBooNE Collaboration, P.~Abratenko \textit{et al.},
  ``Measurement of single charged pion production in charged-current
  $\nu_\mu$-Ar interactions with the MicroBooNE detector,''
  \textit{Phys.\ Rev.\ D} \textbf{113}, 032007 (2026),
  arXiv:2509.03628 [hep-ex].

\bibitem{pt_note} I.~Pophale and J.~Nowak, ``Measurement of proton-tagged charged-current
  single-charged-pion production on argon in the NuMI beam with \uboone{},'' companion analysis
  note (secondary analysis under validation).
\end{thebibliography}

\section{The \uboone{} detector}
\label{sec:detector}

\subsection{Operating Principle}
\label{sec:lartpc}

A liquid argon time projection chamber records three-dimensional ionisation images
of charged-particle tracks by exploiting two properties of ultra-pure liquid argon:
\begin{enumerate}[noitemsep]
  \item Ionisation electrons produced by a traversing charged particle are not
    captured by argon atoms and can drift centimetres or metres under an applied
    electric field without significant recombination, provided the purity is
    sufficient (O$_2$-equivalent $\lesssim 100$~ppt).
  \item Argon scintillates at 127~nm (vacuum ultraviolet), producing $\sim$24\,000
    photons per MeV of deposited energy, enabling nanosecond-scale timing.
\end{enumerate}
An applied electric field drifts the ionisation electrons to a set of sense-wire
planes at the anode, which record three two-dimensional projections of each track.
The third coordinate (along the drift direction) is reconstructed from the known
drift velocity and the timing of the signals relative to the beam-gate start,
combining the two-dimensional wire images into a full 3D reconstruction.

\uboone{} operated from October 2015 to July 2021 at the Fermilab Accelerator
Laboratory, collecting data in both the Booster Neutrino Beam (BNB) and, as a
parasitic experiment, the NuMI beam~\cite{numi_beam}.  It was the world's
largest operating LArTPC until the start of SBND and ICARUS data-taking.

\subsection{Active Volume and Electric Field}
\label{sec:tpc}

The \uboone{} active volume has rectangular geometry:
$256~\mathrm{cm}$ (drift, $x$) $\times$ $233~\mathrm{cm}$ (height, $y$) $\times$
$1037~\mathrm{cm}$ (beam, $z$), containing $\approx 87$~tonnes of liquid argon at
$\approx 87$~K.  The cathode is a stainless-steel plane at $x = 256$~cm held at
$-70$~kV; the anode wire planes are at $x = 0$.  The resulting 273~V/cm drift
field gives an electron drift velocity of $\approx 1.11$~mm/$\mu$s, so the
maximum drift time across the full gap is approximately 2.3~ms.

Liquid argon purity is maintained by continuous recirculation through oxygen and
water absorbers, keeping the electron lifetime well above 5~ms (corresponding to
O$_2$-equivalent contamination below $\sim$100~ppt).  At this purity level,
charge attenuation over the full 2.56~m drift is less than 10\%.

\subsection{Wire Readout and Signal Formation}
\label{sec:wires}

Three wire planes at the anode sample the drifting charge:
\begin{itemize}[noitemsep]
  \item \textbf{U plane} (first induction): 2400 wires at $+60^\circ$ from the
    vertical, 3~mm pitch.
  \item \textbf{V plane} (second induction): 2400 wires at $-60^\circ$ from the
    vertical, 3~mm pitch.
  \item \textbf{Y plane} (collection): 3456 vertical wires, 3~mm pitch.
\end{itemize}
Induction-plane wires are biased to be transparent to the drifting electrons,
producing bipolar (derivative-like) induced signals; collection-plane wires absorb
the charge and produce unipolar signals.  The combination of two stereo induction
views at $\pm60^\circ$ and a vertical collection view provides sufficient 2D
information in each plane to resolve ambiguities during 3D reconstruction.

Front-end electronics (application-specific integrated circuits, ASICs) are
mounted inside the cryostat on the wire planes and cold-multiply each channel at
$\approx 87$~K, reducing thermal noise.  The analogue signals are digitised at
16~MHz inside the cryostat by 12-bit ADCs and transmitted over $\sim$110~m of
cold cables to warm readout electronics.  The data stream is sampled into a
circular buffer; beam triggers and cosmic triggers select 4.8~ms readout windows
(corresponding to $\approx$two full drift windows) for permanent storage.

A dead-wire region at $z \in [675.1, 775.1]$~cm interrupts the collection plane
over a 100~cm span due to a shorted wire bundle.  Tracks crossing this region
have degraded hit-counting and range-based momentum estimation; the fiducial
volume (Section~\ref{N-sec:fv}) excludes this region.

\subsection{Photon Detection System}
\label{sec:pmt}

Thirty-two 8-inch Hamamatsu R5912-MOD photomultiplier tubes (PMTs) are installed
behind the collection-plane wire frame facing the active volume.  The PMTs detect
the 127~nm argon scintillation light via a tetraphenyl-butadiene (TPB)
wavelength-shifting coating that converts VUV photons to visible ($\sim$430~nm)
light.  The system serves two purposes in this analysis:
\begin{enumerate}[noitemsep]
  \item \textbf{Beam-gate synchronisation}: the PMT flash time identifies which
    1.6~$\mu$s NuMI spill window produced the neutrino interaction.
  \item \textbf{Flash-matching for cosmic rejection}: the Pandora reconstruction
    associates PFParticle hypotheses with the PMT flash closest in time and
    position, suppressing out-of-time cosmic-ray muons that overlap with
    beam-induced events in the long 2.3~ms drift window.
\end{enumerate}

\subsection{Calibration and Corrections}
\label{sec:calibration}

Three independent calibration inputs are used:
\begin{itemize}[noitemsep]
  \item \textbf{UV laser system}~\cite{uboone_laser}: two pulsed laser beams produce
    straight ionisation tracks at precisely known locations.  Deviations of the
    reconstructed track positions from the known trajectories are used to build
    a three-dimensional space-charge distortion map~\cite{sce}, which is applied to all
    reconstructed hits as a position correction.  Distortions reach $\sim$10~cm
    near the anode and cathode edges.
  \item \textbf{Stopping cosmic-ray muons}: the Bragg peak of cosmic muons
    stopping inside the active volume calibrates the absolute energy scale and
    the effective recombination factor as a function of local electric field and
    drift distance.  Through-going muons calibrate relative wire-by-wire gains
    and the electron lifetime.
  \item \textbf{Michel electrons}: $\mu^- \to e^-\bar{\nu}_e\nu_\mu$ decay
    products from stopped muons provide a check of the calorimetric
    reconstruction at $E_e \lesssim 55$~MeV.
\end{itemize}

\subsection{Cosmic Ray Environment and Reconstruction}
\label{sec:cosmics}

Operating at the surface with minimal overburden, \uboone{} is exposed to a
cosmic-ray muon rate of $\approx 5$~kHz through the active volume.  Each 4.8~ms
readout window contains on average $\sim$10 cosmic-ray tracks unrelated to the
beam interaction.  The 1.6~$\mu$s NuMI spill itself contains $\lesssim 0.1$
expected cosmics on average, but the drift-window pile-up is the dominant
background at early selection stages (Section~\ref{N-sec:selection}).

Cosmic rejection proceeds in stages:
\begin{enumerate}[noitemsep]
  \item Flash-matching (PMT timing): candidate neutrino slices must have a
    consistent optical flash within the beam window.
  \item Topological score: the Pandora-based topological neutrino score
    discriminates between neutrino-like vertex structures and
    through-going/stopping cosmic topologies.
  \item Dedicated BDT scores: the MIP and cosmic
    scores trained on track-level and event-level features further suppress
    residual cosmics that passed the topological cut.
\end{enumerate}
The CRT (Cosmic Ray Tagger) sub-detector~\cite{uboone_crt}, consisting of
scintillator panels surrounding the cryostat, provides an additional cosmic
veto but is not used as a cut in the present analysis: the CRT branches are carried
through the event processing as pass-through quantities only, they are populated only
for the CRT-era data (Runs 3--5) and zero for the pre-CRT Run-1 exposure, and a cut would
require a data/MC tagging systematic; the residual beam-off background is subtracted
instead (\noteref{Sec.}{sec:mc}).

Reconstruction uses the Pandora~\cite{pandora,uboone_pandora} multi-algorithm pattern recognition
framework, which groups hits into PFParticles organised in a hierarchy with the
neutrino interaction as the top-level particle.  Track candidates are fitted with
a Kalman-filter-based track fitter to obtain momentum from track curvature
(negligible for $<$300~MeV/$c$ muons at 273~V/cm) and from range for
contained tracks.  Log-likelihood ratio particle identification (LLR PID)~\cite{uboone_llr} using
the $dE/dx$ profile along the track is the primary tool for muon/pion/proton
separation (Section~\ref{N-sec:selection}).

\end{document}
